\documentclass[11pt,a4paper]{article}

% ===== XeLaTeX + Korean =====
\usepackage{kotex}
\usepackage{fontspec}
\setmainfont{Times New Roman}
\setmainhangulfont{AppleMyungjo}
\setsanshangulfont{Apple SD Gothic Neo}

% ===== Layout =====
\usepackage[a4paper,top=18mm,bottom=18mm,left=20mm,right=20mm]{geometry}
\usepackage{fancyhdr}
\setlength{\headheight}{24pt}
\usepackage{enumitem}
\usepackage{mdframed}
\usepackage{array}

\setlength{\parindent}{1em}
\linespread{1.18}

% ===== Header / Footer =====
\pagestyle{fancy}
\fancyhf{}
\renewcommand{\headrulewidth}{0.6pt}
\fancyhead[L]{\sffamily\bfseries 2026 고1 영어 변형 — 지문 29}
\fancyhead[R]{\sffamily 문학의 힘}
\renewcommand{\footrulewidth}{0pt}
\fancyfoot[C]{\framebox[16mm][c]{\thepage}}

% ===== helpers =====
\newcommand{\qno}[1]{\par\vspace{8pt}\noindent{\sffamily\bfseries #1.}\ }
\newcommand{\instr}[1]{{\sffamily #1}\par\vspace{3pt}}
\newlist{choices}{enumerate}{1}
\setlist[choices]{label=,leftmargin=1.5em,itemsep=2pt,topsep=3pt,parsep=0pt}

\newmdenv[linewidth=0.6pt,roundcorner=3pt,innertopmargin=7pt,innerbottommargin=7pt,
          innerleftmargin=9pt,innerrightmargin=9pt,skipabove=5pt,skipbelow=5pt]{pbox}
\newcommand{\nemo}[1]{\fbox{\,#1\,}}

\begin{document}

\begin{center}
{\fontsize{15}{17}\selectfont\sffamily\bfseries [지문 29] 다음 글을 읽고 물음에 답하시오.}
\end{center}
\vspace{4pt}

% ===== 공통 지문 (네모 + 빈칸 + 밑줄 동시 삽입) =====
\begin{pbox}
Literature, in its essence, is a testament* to the enduring power of human creativity and expression. It (A)\,\nemo{transcends / transcending} time and space, connecting us with the thoughts and emotions of people from all walks of life, across centuries and cultures. (B)\,\nemo{Despite / Although} the rise of new technologies and the ever-changing landscape of media, literature has (C)\,\nemo{retained / been retaining} its relevance and significance as a timeless art form. One reason for literature's enduring power is its ability to capture and convey \rule{5cm}{0.4pt}. Through the written word, authors \underline{having} the unique ability to explore the depths of human emotion, the nuances*** of relationships, and the challenges and triumphs of the human spirit. Literature allows us to see ourselves reflected in the pages of a book, to empathize with characters different from ourselves, and to gain a deeper understanding of the human condition.

\vspace{4pt}
\noindent\small * testament: 증거 \quad ** transcend: 초월하다 \quad *** nuance: 미묘한 차이
\end{pbox}

% ===================== Q1 =====================
\qno{1}\instr{(A), (B), (C)의 각 네모 안에서 어법에 맞는 표현으로 가장 적절한 것은?}
\begin{center}
\renewcommand{\arraystretch}{1.25}
\begin{tabular}{c c c c}
 & (A) & (B) & (C) \\
① & transcends & Despite & retained \\
② & transcends & Despite & been retaining \\
③ & transcends & Although & retained \\
④ & transcending & Despite & retained \\
⑤ & transcending & Although & been retaining \\
\end{tabular}
\end{center}

% ===================== Q2 =====================
\qno{2}\instr{다음 빈칸에 들어갈 말로 가장 적절한 것은?}
\begin{choices}
\item ① the historical events of past civilizations
\item ② the technical skills required for creative writing
\item ③ the intricacies of what it means to be human
\item ④ the entertainment needs of a modern audience
\item ⑤ the latest scientific discoveries of our time
\end{choices}

% ===================== Q3 (서술형) =====================
\qno{3}\instr{[서술형] 윗글에서 어법상 틀린 부분을 찾아 올바르게 고쳐 쓰시오. (밑줄 친 부분 기준)}

\vspace{4pt}
\noindent $\Rightarrow$ \rule{\linewidth}{0.4pt}

% ===================== Q4 =====================
\qno{4}\instr{윗글의 제목으로 가장 적절한 것은?}
\begin{choices}
\item ① Literature: A Timeless Mirror of Human Experience
\item ② How New Technologies Have Replaced Written Art
\item ③ Why Historical Accuracy Matters More Than Emotion
\item ④ Literature as Entertainment for Modern Audiences Only
\item ⑤ The Decline of Reading Across Centuries and Cultures
\end{choices}

\end{document}
